% ###################
% ## Filename: sCM_drivers.tex
% ## Location: D:\Project\Prej_Teaching\ISE\SupplyChainMgt\presentation\scm_02_drivers
% ## Created: Dec 14, 2010
% ## Last updated: Dec 14, 2010
% ## Figure Location: ..\Figures\
% ## Note: Presentation, Teaching
% ###################


%\documentclass[10pt,presentation]{beamer}
%\documentclass[10pt,handout]{beamer}
\documentclass[10pt,handout,notes=show]{beamer}

%============================================================================================================--
\usepackage{fontspec}
\usepackage{xunicode}
\usepackage{xltxtra}
\usepackage{array}
\usepackage{xcolor}
\usepackage{verbatim}
\usepackage{beamerthemesplit}
\usepackage{beamerbasetemplates}
\usepackage{graphics}
\usepackage{graphicx}
\usepackage{hyperref}
\usepackage{multirow}
\usepackage{tikz,pgfplots}
\usepackage{algorithm}
\usepackage{algorithmic}
\usepackage{soul,xcolor}
\usepackage{colortbl}
\usepackage{multicol}

\XeTeXlinebreaklocale"th"
\defaultfontfeatures{Scale=1.5, Mapping=tex-text} %{Scale=1.23}



\usetheme{Madrid}  % JuanLesPins  Rochester
           %  Berkeley  Palo Alto    with sidebar and top
           % Goettingen  Marburg     with sidebar
           %  Copenhagen  Luebeck  Warsaw

\setbeamertemplate{headline}
{%
    \begin{beamercolorbox}{section in head/foot}
        \vskip-2pt \hskip5pt
        \insertsectionnavigationhorizontal{\textwidth}{}{}
        \hskip5pt \vskip-2pt
    \end{beamercolorbox}
}%
\setbeamertemplate{footline}
{%
    \begin{beamercolorbox}{section in head/foot}
        \vskip-2pt \hskip5pt
        \insertshorttitle \hfill \insertshortdate \hskip10pt \insertframenumber / \inserttotalframenumber
        \vskip-2pt\hskip5pt
    \end{beamercolorbox}
}%
\definecolor{alertedOrange}{RGB}{255,36,0}
\setbeamercolor{alerted text}{fg=alertedOrange}
\useinnertheme{rounded}
\usecolortheme{rose}
%\setbeamercovered{transparent}
\usefonttheme{structuresmallcapsserif}
\definecolor{LRed}{rgb} {1,.8,.8}
\definecolor{LBlue}{rgb}{.8,.8,1}

%%%%%%%%%%%%%%%%%%%%%%%%%%%%%%%%%%%%%%%%%%%%%%%%%%%%%%%%%%%%%%%%%%%%%%%%%%%%%%%
% Notation:
\newtheorem{observation}{Observation}
\newcommand{\figureLocation}[1]{../figures/#1}
\newcommand{\figureLoc}[1]{../figures/#1}
\newcommand{\refLocation}[1]{../../../../Reference/#1}
\newcommand{\NextLine}{\vskip 0.05\textheight}
\newcommand{\NextHalfLine}{\vskip 0.02\textheight}
\newcommand{\RemoveLine}{\vskip -0.1\textheight}
\newcommand{\RemoveHalfLine}{\vskip -0.02\textheight}
\newcommand{\BackTab}{\hskip -0.15\textwidth}
\newcommand{\HalfBackTab}{\hskip -0.11\textwidth}
\newcommand{\DualFigureSize}{1.10\textwidth}
\newcommand{\MiniPageSize}{0.52\textwidth}
\newcommand{\code}[1]{\textcolor{blue}{#1} }
\newcommand{\rightNote}[2]{\hfill \uncover<#1>{\textbf{\code{\#\# {#2}  }}}}
\newcolumntype{P}[1]{>{\raggedright\arraybackslash}p{#1}<{}}

\author[Kittithreerapronchai]{%
  Oran~Kittithreerapronchai\inst{1} }
\institute[IE Chula]{
  \inst{1}%
  Department of Industrial Engineering, Chulalongkorn University\\
  Bangkok 10330 THAILAND \\
  }
\subject{}
\date{}

\setbeamertemplate{enumerate items}[default]
\setbeamertemplate{bibliography item}[text]
\bibliographystyle{alpha}
\newcommand{\textFont}[1]{\footnotesize{#1}}
\setcounter{MaxMatrixCols}{20}

%%%--------------------------------

\makeatletter
\defbeamertemplate{note page}{infolines}
{%
  {%
    \small
    \insertvrule{.25\paperheight}{white!90!black}
    \vskip-.25\paperheight
    \nointerlineskip
    \vbox{
      \hfill\insertslideintonotes{0.25}\hskip-\Gm@rmargin\hskip0pt%
      \vskip-0.25\paperheight%
      \nointerlineskip
      \begin{pgfpicture}{0cm}{0cm}{0cm}{0cm}
        \begin{pgflowlevelscope}{\pgftransformrotate{90}}
          {\pgftransformshift{\pgfpoint{-2cm}{0.2cm}}%
          \pgftext[base,left]{\footnotesize\the\year-\ifnum\month<10\relax0\fi\the\month-\ifnum\day<10\relax0\fi\the\day}}
        \end{pgflowlevelscope}
      \end{pgfpicture}}
    \nointerlineskip
    \vbox to .25\paperheight{\vskip0.5em
      \hbox{\insertshorttitle[width=8cm]}%
      \setbox\beamer@tempbox=\hbox{\insertsection}%
      \hbox{\ifdim\wd\beamer@tempbox>1pt{\hskip4pt\raise3pt\hbox{\vrule
            width0.4pt height7pt\vrule width 9pt
            height0.4pt}}\hskip1pt\hbox{\begin{minipage}[t]{7.5cm}\def\breakhere{}\insertsection\end{minipage}}\fi%
      }%
      \setbox\beamer@tempbox=\hbox{\insertsubsection}%
      \hbox{\ifdim\wd\beamer@tempbox>1pt{\hskip17.4pt\raise3pt\hbox{\vrule
            width0.4pt height7pt\vrule width 9pt
            height0.4pt}}\hskip1pt\hbox{\begin{minipage}[t]{7.5cm}\def\breakhere{}\insertsubsection\end{minipage}}\fi%
      }%
      \setbox\beamer@tempbox=\hbox{\insertshortframetitle}%
      \hbox{\ifdim\wd\beamer@tempbox>1pt{\hskip30.8pt\raise3pt\hbox{\vrule
            width0.4pt height7pt\vrule width 9pt
            height0.4pt}}\hskip1pt\hbox{\insertshortframetitle[width=7cm]}\fi%
      }%
      \vfil}%
  }%
  \vskip.25em
  \nointerlineskip
  %\setmainfont{TH Sarabun New}
  \begin{minipage}{\textwidth} % this is an addition
  \footnotesize{\fontspec{TH Sarabun New} \insertnote}
  \end{minipage}               % this is an addition
  %\setmainfont{Computer Modern}

}
\makeatother

\setbeamertemplate{note page}[infolines]

\title[LSCM v3.0]{Lecture 02:\\ Supply Chain Perspective and Strategy}
\hypersetup{
  pdftitle={02 Supply Chain Perspective and Strategy},
  pdfsubject={2104509lscm & 2184408scm handout},
  pdfproducer={XeLatex/WinEdt},
  pdfcreator={oran.k@chula.ac.th}
}


\begin{document}

%==============================================================================================================
\begin{frame}[t]
  \titlepage \label{Slide:Title}

  \NextLine   \NextLine
  \begin{flushright}
    \tiny{ \textbf{last updated:} \today } \hskip3ex
  \end{flushright}

  \note{
    \textbf{Purpose:} Big picture of SCM (business perspective) and important strategy\\
    \textbf{Leading Question:} What your BBA friends know about SCM?\\
    \begin{itemize}
      \item \textbf{Re-Cap:} SCM is around us and dynamic
      \item \textbf{Diff btw Logistic \& SCM:} key diff including strategy, manufacturing, connectivity
      \item \textbf{Disclaimer:} strategy is NOT OKC forte; it also takes time to review; no right or wrong; no middle road
    \end{itemize}
  }
\end{frame}

%==============================================================================================================
\begin{frame}[t]
  \frametitle{Outline}

  \tableofcontents

  \vfill
  \begin{flushright}
    \tiny{\textbf{General Reference:} \cite{jacobs2010operations}~\cite{ballou2007evolution} \cite{chopra2007supply} \cite{goetschalckx2011supply} } \hskip3ex
  \end{flushright}

  \note{
    \textbf{Purpose:} Outline topics\\
    \textbf{Leading Question:} <<none>>\\
    \begin{itemize}
      \item \textbf{Perspective:} how to view about SCM; flow, process, and SCOR
      \item \textbf{SCOR:} a managerial tool to implement SCM in company
      \item \textbf{SCM Strategy:} how SCM link to business strategy
      \item \textbf{SCM Fit:} how to craft a good SCM strategy what we should think
      \item \textbf{SCM Matrix:} how to measure and improve
      \item \textbf{Bullwhip:} phenomenal that unique to SCM, not logistics
      \item \textbf{\code{EXTRA}: How improve your SCM:} how to think in terms of short-long terms and design-execute
    \end{itemize}
  }
\end{frame}

%============================================================================================================--
\section[Perspective]{Perspectives in Supply Chain Management}
%================================================================================================-

\begin{frame}
  \frametitle{SCM Development Time Line}

   \begin{center}
    \includegraphics[width=0.85\linewidth]{\figureLocation{scm_strategy/devlopSCM.png} }
   \end{center}

  \begin{flushright}
    \tiny{\textbf{source:} Mamun Habib. ``Supply Chain Management-Applications and Simulations''} \hskip10em
  \end{flushright}

  \note{
    \textbf{Purpose:} Evaluation of SCM\\
    \textbf{Leading Question:} How our SCM perspective change company view?\\
    \begin{itemize}
      \item \textbf{traffic:} supp $\rightarrow$ company $\rightarrow$ cust:  manual/ many carries
      \item \textbf{logistic:} suppOfsupp $\rightarrow~\cdots$ company $\cdots~\rightarrow$ custOfcust
      \item[] max profit/low cost + measure / few carriers/ bar code
      \item \textbf{supply chain:} suppOfsupp $\rightarrow~\cdots$ company $\cdots~\rightarrow$ custOfcust
      \item[] value (speed/scale/outsoucing ) WMS/
      \item \textbf{collaboration:} how to spring new business (data)
    \end{itemize}
  }
\end{frame}

%============================================================================================================--

\begin{frame}
  \frametitle{Changing view of company}

  \begin{center}
  \begin{tabular}{lr}
    \includegraphics [width=0.35\textwidth]{\figureLocation{scm_strategy/tradView.png}}~
    &
    ~\includegraphics [width=0.35\textwidth]{\figureLocation{scm_strategy/scmView.png}}
  \end{tabular}
  \end{center}

  \begin{center} \small{
    \begin{tabular}{l|rr} \hline
    \multicolumn{1}{c}{}
                     & \multicolumn{1}{c}{\textbf{Traditional}}
                     & \multicolumn{1}{c}{\textbf{Value creation}} \\ \hline
    \textbf{Parties} & \alert{one company}      & \alert{every companies} \\
    \textbf{Concept} & winner takes all         & focus on \alert{my 'things'} \\
    \textbf{Lv Decision} & operation            & \alert{strategic} \\
    \textbf{Values}  & cost                     & \alert{revenue stream}\\[2pt]
                     &                           & , engagement, \alert{equity} \\[5pt]
    \textbf{Tools}   & price \& quality         & speed \& information \\ \hline
    \end{tabular} }
  \end{center}

  \NextHalfLine
  \begin{block}{Question 1}{}
    Do your customers care? If not, why bother?\\
    \uncover<2-|handout:1>{\hfill \code{Effect:} cost-availability, reliability-flexibility, quality, convenience, fluctuation \hfill}
  \end{block}

  \note{
    \textbf{Purpose:} Evaluation of SCM\\
    \textbf{Leading Question:} How our 'flow' perspective change company view?\\
    \begin{itemize}
      \item \textbf{Impacts/Stories}
        \begin{itemize} \scriptsize{
          \item Grocery industry could save \$30 billion  by using effective supply chain strategies
          \item Laura Ashley turned its inventory 10 times, five times faster than 3 years ago
          \item In 1995, Compaq lost \$.5 to \$1 billion in sales because of laptop misplacing
          \item When AMD introduced 1 gig processor, the 800 mb price dropped by 30\%
          \item P\&G saved $\approx$ \$65 million by collaboration (better supply and demand matching)
          \item Dell received \$2000 from a customer for a computer
          \item In 1998, Blockbuster increased 5\% of the rental market at the start by sharing revenue with DVD companies
          \item 7-Eleven invested in IT to manage supplies and demands
          \item Toyota's plant could shift auto manufacturing for domestic market into export aborad }
      \end{itemize}
      \item \textbf{Trad \textsc{VS} Value:} what you must aware yourself
      \item \textbf{Consumer = King:} important of experience, e.g. Amazon, Starbucks,
    \end{itemize}

  }
\end{frame}

%================================================================================================-
\begin{frame}
  \frametitle{For Customer, Supply Chain is}

  Goals of Supply Chain is to ensure that \alert{customers} got
  \begin{itemize}
    \item[] the right \alert{item}         (national \textsc{VS} store brands)
    \item[] in the right \alert{quantity}  (EA \textsc{VS} DOZEN)
    \item[] at the right \alert{price}     (7Eleven \textsc{VS} Makro)
    \item[] at the right \alert{place}     (ENG canteen \textsc{VS} Airport)
    \item[] at the right \alert{time}      (iPhone Camping)
    \item[] in the right \alert{condition} (edible gelatine \textsc{VS} photo gelatine)
    \item[] at the optimum cost to the \alert{organization(s)}
  \end{itemize}

  \begin{flushright}
    \tiny{\textbf{source:} CSCMP, ``SCM Goals –- The Seven Rights'' 2011 \hskip10em}\\
  \end{flushright}

  \note{
    \textbf{Purpose:} Introduce to SCM from customer's view\\
    \textbf{Leading Question:} What does a customer care?\\
    \begin{itemize}
      \item \textbf{why use 'Right':} why use 'right' $\rightarrow$ at what consumer perceptive
      \item \textbf{Right Product/Condition/Price:} customization, complexity, unknown needs, e.g.  ergoChair, iPad
      \item \textbf{Right Qty:} key for negotiation $\rightarrow$ \alert{away enough}
      \item \textbf{Right Place:} global demand, global supply $\rightarrow$ \alert{away now}
      \item \textbf{Right Time:} JIT, shorten LT, $\rightarrow$ \alert{away now}
      \item \textbf{SCM \textsc{VS} Acct Time measurement:} used to be day; now hour
      \item \textbf{Organization:} target customers; business strategy
    \end{itemize}
    %% (Chain 1st Quality, Intel, Harddrive)
  }
\end{frame}

%============================================================================================================--
\begin{frame}
  \frametitle{Other perspective of SCM}

    \begin{description}
      \item[Flows:] movement of 'things' along SC (traditional views)
      \begin{minipage}{0.8\textwidth} \small{
        \begin{tabular}{l|rr} \hline
                                    & Upstream  & Downstream \\ \hline
        \textbf{Materials:}   & return, repair, disposal&  RM, WIP, FG   \\
        \textbf{Information:} & sales, promotion plan   &  capacity, delivery schedule \\
        \textbf{Funds/Values:}& payment, consignment    & credit, invoice, payment\\ \hline
        \end{tabular} }
      \end{minipage} \NextHalfLine
      \item[Process:] supply chain activity repeats itself (system views)
      \begin{itemize}
        \item \textbf{Cycle order:} process in SC as series of cycles each performed by interfaced parties
        \item \textbf{Push/Pull:} processes in SC as responding (\alert{pull}) or in anticipating (\alert{push}) to order \NextHalfLine
      \end{itemize}
      \item[SCOR:] supply chain activity can be grouped by business function
       \begin{itemize}
        \item standardizing \& combining flow and process
        \item introducing KPIs
        \item allowing benchmark between companies
       \end{itemize}
    \end{description}

  \note{
    \textbf{Purpose:} Introduce three related perspective\\
    \textbf{Leading Question:} Why 'flow' perspective of SCM was sexy and new? How those who worked in the area never though about it?\\
    \begin{itemize}
      \item \textbf{b/g:} limited info \& ecosystem $\rightarrow$ focus on a single company \& single department; logistics $\rightarrow$ cost center
      \item \textbf{Flow:} visualizing how fund/info/physical connecting between parties $\rightarrow$ who create \alert{value}
      \item \textbf{Process:} routining operations based on mistake or experiance, cycle of XXX (work on the same procedure without improvement) $\rightarrow$ how to \alert{done}
      \item \textbf{SCOR:} flow + process + furthermore
    \end{itemize}
  }
\end{frame}

%============================================================================================================
\usetikzlibrary{arrows,positioning,decorations.pathmorphing,calc,shapes}
\renewcommand{\textFont}[1]{\footnotesize{#1}}
%================================================================================================-
\begin{frame}
  \frametitle{Cycle or Process Perspective }

  \begin{center}
  \begin{tabular}{lc}
    \includegraphics[width=0.50\textwidth]{\figureLocation{scm_strategy/functionLogistics.png}}
    &
    \includegraphics[width=0.35\textwidth]{\figureLocation{scm_strategy/CycleProcess.png}}
    \\
    \tiny{~\hskip10em \textbf{source:} Martin Christopher, 2011~\cite{christopher2016logistics} }
    &
    \tiny{~\hskip10em \textbf{source:} Chopra and Meindl 2010~\cite{chopra2007supply}\hskip0em }
  \end{tabular}
  \end{center}

  \begin{block}{Question 2}{}
    What are triggers/information of these processes?
  \end{block}

  \note{
    \textbf{Purpose:} Introduce to easy business process\\
    \textbf{Leading Question:} How SCM operators work?\\
    \begin{itemize}
      \item \textbf{Cycle (right figure):} synchronize and repeat cycles of customer order/ replenishment/ manufactureing/ purchasing
      \item \textbf{Note on Cycles:} lean \& mean $\rightarrow$ short cycle $\rightarrow$ new innovation, high inventory turnover
      \item \textbf{Typical Process (left figure):} physical flow is easy. triggers of these processes are information.
      \item \textbf{Question:} What are they? What do they derive from?
      \item \textbf{Important Logistics Terms:} Order, GI, GR, Invoice, PO, Bill of Lading, Manifest, Freight Bill, L/C
    \end{itemize}
  }
\end{frame}

%=================================================================================================
\begin{frame}
  \frametitle{Production/Inventory Systems}

  \uncover<2->{
  \begin{center}
  \begin{tabular}{rll} \hline
    ~         & Push                    & Pull \\ \hline
    Shipment  & large \& infrequent     & small \& frequent \\
    Suppliers & many \& short term      & few \& long term \\
    Information& minimal exchange       & constant exchange\\
    Geographic& unimportant             & adjacent \\
    Application& MRP & JIT \\ \hline
  \end{tabular}
  \end{center}
  }

  \NextLine
  \textbf{Assembly Line Simulation}

  \note{
    \textbf{Purpose:} Understand Pull and Push and Push-Pull Boundary\\
    \textbf{Leading Question:} <<none>> \\
    \begin{itemize}
      \item push \textsc{VS} pull: how to estimate demand"
      \item \textbf{push:} long LT, low value \& large quality, common $\rightarrow$ high inventory
      \item \textbf{pull:} quick response, high value, customize, spec $\rightarrow$ individual
      \item \textbf{Simulation:} where is 'push-pull' boundary
    \end{itemize}
  }
\end{frame}


%================================================================================================-
\begin{frame}
  \frametitle{Push/Pull Perspective }

  \begin{center}
  \begin{tabular}{lr}
    \only<1|handout:1>{\includegraphics[height=0.25\textheight]{\figureLocation{scm_strategy/PushPullProcess.png}}}
    \only<2-|handout:0>{\includegraphics [height=0.25\textheight]{\figureLocation{scm_strategy/pushpullStr.png}}}
    &
    \only<1-|handout:1>{\includegraphics [height=0.25\textheight]{\figureLocation{scm_strategy/demandVSforecast.png}}}
    \\
    \tiny{~\hskip10em \textbf{source:} Chopra and Meindl 2010~\cite{chopra2007supply}\hskip0em }
    &
    \tiny{~\hskip10em \textbf{source:} Martin Christopher, 2011~\cite{christopher2016logistics} }
  \end{tabular}
  \end{center}

  \NextHalfLine
  \begin{block}{Question 3: Thinking Points}{}
  What are good characteristics of push (pull) products? and How to leverage?

  \end{block}

  \note{
    \textbf{Purpose:} diff strategy for diff business/product\\
    \textbf{Leading Question:} what is the right strategy for me?\\
    \begin{tabular}{cc}
      \includegraphics[width=0.4\textwidth]{\figureLocation{scm_strategy/PushPullProcess.png}} &
      \includegraphics[width=0.4\textwidth]{\figureLocation{scm_strategy/pushpullStr.png}} \\
    \end{tabular}
    \begin{itemize}
      \item Upstream $\rightarrow$ Push (far from end customer, low value, move in bulk)
      \item Downstream $\rightarrow$ Pull (near end customer, high value)
    \end{itemize}
    \textbf{How should we think?}
    \begin{itemize}
      \item point of used (mfg \textsc{VS} spar part in auto industry)
      \item leadtime \& rate of consumption (takt time)
      \item qty / lifecycle of items
    \end{itemize}
  }
\end{frame}

%================================================================================================-
\begin{frame}
  \frametitle{Uncertainty and Scales }

  \begin{center}
    \includegraphics [width=0.6\textwidth]{\figureLocation{scm_strategy/pushpullMat.png}}
  \end{center}
  \begin{flushright}
    \tiny{\textbf{source:} Martin Christopher, 2011~\cite{christopher2016logistics} \hskip10em}
  \end{flushright}

  \begin{itemize}
    \item \textbf{Pull Strategy (I):} high demand uncertainty without economy of scale
    \item \textbf{Push Strategy (III):} low demand uncertainty with economy of scale
  \end{itemize}

  %%%------
  \note{
    \textbf{Purpose:} Matrix btw uncertainty and scale\\
    \textbf{Leading Question:} Any guideline to mgt push/pull strategy\\
    \begin{itemize}
      \item \textbf{Push =} Supply Chain Planning processes are applied $\rightarrow$ ok to 'guess'
      \item \textbf{Push:} low uncertainty, service level not an issue, focus on cost minimization, long lead times
      \item[] \hskip3ex $\bullet$ better utilizing resources
      \item[] \hskip3ex $\bullet$ minimizing inventory, transportation, and production costs
      \item \textbf{Pull =} Order-fulfillment processes are applied $\rightarrow$ \alert{not ok} to 'guess'
      \item \textbf{Push:} $<<$\emph{opposite to push}$>>$
      \item[] \hskip3ex $\bullet$ quality matter; error can propagate (know first)
      \item[] \hskip3ex $\bullet$ minimizing risk; induce experiance
    \end{itemize}
  }
\end{frame}


%============================================================================================================--
\section[SCOR]{Supply Chain Operation Reference Model}
%==============================================================================================================

\begin{frame}
  \frametitle{Overview of SCOR model}

  \begin{center}
    \only<1|handout:0>{\includegraphics [height=0.5\textheight]{\figureLocation{scm_strategy/scor_model_concept01.png}}}
    \only<2-|handout:1>{\includegraphics[height=0.5\textheight]{\figureLocation{scm_strategy/scor_model_cV11.png}} }
  \end{center}
  \begin{flushright}
    \tiny{\textbf{source:}  Cohen, S. and Roussel, J 2001. ``Strategic Supply Chain Management'' \hskip5em}\\
  \end{flushright}

  \begin{itemize}
    \item \textbf{What:} hierarchy business process framework
    \item \textbf{Scopes:} all \alert{customer's interaction}, \alert{product transactions}, \alert{market interactions}
    \item \textbf{Benefit:} std process, best practice, matrices for benchmark,
    \item \textbf{Issues:} no details operation, ignored sale and markets, R\&D
  \end{itemize}

  \note{
    \textbf{Purpose:} Introduce to SCOR model\\
    \textbf{Leading Question:} what are minimal requirement for SCM student to know\\
    \begin{tabular}{cc}
      \includegraphics [width=0.4\textwidth]{\figureLocation{scm_strategy/scor_model_concept.png}} &
      \includegraphics [width=0.4\textwidth]{\figureLocation{scm_strategy/SCOR_model_bV11.png}} \\
    \end{tabular}
    \begin{itemize}
      \item \textbf{SCOR:} Supply Chain and Operation Reference, first to address the whole company with business function, common term,
      \item \textbf{4 Levels:} processes, subprocess, activities, operation process (skipped)
      \item \textbf{Two levels:} present a broad scope/focus on where business are common
      \item \textbf{Make:} the hardest part, typically ignore,
      \item \textbf{Serve as:} check point and framework
      \item \textbf{Bad for:} starting point (too comprehensive model from beginning to the end)
      \item \textbf{Goal:}  To develop a standard supply-chain process reference model enabling effective communication
      \item \textbf{Assume}  Training, Quality, IT, admin
      \item \textbf{Level:} Process, Process Element, Task, Activity (implement)
    \end{itemize}
  }

\end{frame}

%==============================================================================================================
\begin{frame}
  \frametitle{SCOR model: Semi-Perfect Concept }

  \begin{center}
    \only<1|handout:0>{\includegraphics [height=0.5\textheight]{\figureLocation{scm_strategy/scm_framework.png}}}
    \only<2|handout:1>{\includegraphics [height=0.5\textheight]{\figureLocation{scm_strategy/scor_lv0.png}}}
    \only<3|handout:0>{\includegraphics [height=0.5\textheight]{\figureLocation{scm_strategy/scor_lv1.png}}}
    \only<4|handout:0>{\includegraphics [height=0.5\textheight]{\figureLocation{scm_strategy/scor_lv2.png}}}
  \end{center}
  \begin{flushright}
    \tiny{\textbf{source:}  Cohen, S. and Roussel, J 2001. ``Strategic Supply Chain Management'' \hskip5em}\\
  \end{flushright}


  \begin{itemize}
    \item \textbf{Lv1: Model} set scope, content, \alert{performance}, product-segment
    \item \textbf{Lv2: Strategy} decide configuration, i.e., strategy, capability
    \item \textbf{Lv3: Process} identify key business process, i.e., input, output, \alert{best practice}
    \item \textbf{Lv4: Activity} improve activates using lean, TQM, \alert{benchmark}
  \end{itemize}


  \note{
    \begin{tabular}{cc}
    \includegraphics[width=0.4\textwidth]{\figureLocation{scm_strategy/scm_framework.png}} & \includegraphics[width=0.4\textwidth]{\figureLocation{scm_strategy/scor_lv.png}}  \\
    \end{tabular}
    \textbf{Purpose:} Introduce to SCOR model\\
    \textbf{Leading Question:} why SCOR model is perfect but incompleted\\

    \begin{itemize}
      \item \textbf{Hierarchy:} connected level (process and KPIs)
      \item \textbf{Focus:} Lv1 $\rightarrow$ key process, value proposition, focus
      \item \textbf{Config:} Lv2 $\rightarrow$ main business?
      \item \textbf{Activity:} Lv3 $\rightarrow$ how to process, input, output, data
    \end{itemize}
  }

\end{frame}

%==============================================================================================================
\begin{frame}
  \frametitle{SCOR model Lv 1 Metric}

  \begin{center}
    \includegraphics [width=0.70\textwidth]{\figureLocation{scm_strategy/scor_matrixLv1.png}} \\
  \end{center}
  \begin{flushright}
    \tiny{\textbf{source:} \href{http://supply-chain.org/scor}{\code{http://supply-chain.org/scor}} \hskip5em}\\
  \end{flushright}

  %%%----
  \note{
    \textbf{Purpose:} select matrix and KPI\\
    \textbf{Leading Question:} what is the main result of Lv1\\
    Level 1 defines scope of process
    \textsc{Customer-Facing}
    \begin{itemize}
      \item \textbf{Reliability:} fill rate, perfect order
      \item \textbf{Responsiveness:} fulfilment lead time
      \item \textbf{Flexibility:} supply chain response time
    \end{itemize}

    \NextLine
    \textsc{Internal-Facing}
    \begin{itemize}
      \item \textbf{Cost:} supply chain cost, value-added, return cost
      \item \textbf{Assets:} cash-to-cash, inventory day of supply,
    \end{itemize}
  }

\end{frame}

%==============================================================================================================
\begin{frame}
  \frametitle{SCOR model Lv 1 Scope}

    \begin{itemize}
    \item \textbf{Plan:} balance aggregate demand and supply $\rightarrow$ actions
    \begin{itemize}
      \item assess \& acquire \alert{resources} within requirements
      \item plan inventory and distribution and communicate plan
    \end{itemize}
    \item \textbf{Source} procure components of goods for planned demands
    \begin{itemize}
      \item \alert{schedule} deliveries, transfer, contract, payment
      \item certify, select \& assess suppliers (except ETO)
    \end{itemize}
    \item \textbf{Make:} transform product to a finished goods to meet demands
    \begin{itemize}
      \item schedule production, resource, \& infrastructure
      \item cover assemble, repair, overhaul, recycle
    \end{itemize}
    \item \textbf{Deliver:} react, maintenance, and fulfillment of orders
    \begin{itemize}
      \item validation, consolidate, pick, pack, ship and invoice the ordered quantity
    \end{itemize}
    \item \textbf{Return:} returning or receiving returned products
    \begin{itemize}
      \item decide to return good \& container
      \item perform return, substitution, refund, restock
    \end{itemize}
    \item \textbf{Enable:} managerial activities
    \begin{itemize}
      \item business rules \& procurement
      \item mgt of performance, resources, data, contracts, network, compliance, risk
    \end{itemize}
  \end{itemize}
  \begin{flushright}
    \tiny{\textbf{source:} \href{http://supply-chain.org/scor}{\code{http://supply-chain.org/scor}} \hskip5em}\\
  \end{flushright}

  \note{
    \textbf{Purpose:} Introduce to Level 1 \\
    \textbf{Leading Question:} <<none>>\\
    Because a company does more than things, it is important to identify how each BU integrates and what need to be doing in a company.
    Level 1 defines scope of process
    \begin{itemize}
      \item \textbf{Plan:} Input/Process/Output planning and management; financial \& SC alignment
      \item \textbf{Source:} Make-to-Stock, Make-to-Order, Assembly-to-Order, Engineering-to-Order
      \item \textbf{Deliver:} order management, transportation management, and distribution management
      \item[] \hskip3ex $\bullet$ \textsl{delivery retail products} processes used to acquire, merchandise, and sell finished goods at a retail store
      \item \textbf{Make:} $<<$\emph{similar to source}$>>$
      \item \textbf{Return:} of what (RM to source or FG from delivery). Types (defect/maintainace/excess) This depends on strategy and nature of business (underwear, period)
      \item \textbf{Enable:} supporting and gauging, measuring, and improving the effectiveness of SCM
    \end{itemize}
  }

\end{frame}

%==============================================================================================================
\begin{frame}
  \frametitle{SCOR model: Level 2 Config}

  \NextLine
  \begin{center}
    \only<1 |handout:1>{ \includegraphics [width=0.70\textwidth]{\figureLocation{scm_strategy/scor_lv0.png}}}
    \only<2 |handout:0>{ \includegraphics [width=0.70\textwidth]{\figureLocation{scm_strategy/scor_cycleTime.jpg}}}
    \only<3 |handout:0>{ \includegraphics [width=0.70\textwidth]{\figureLocation{scm_strategy/scor_cycleTime01.png}}}
  \end{center}
  \begin{flushright}
    \tiny{\textbf{source:} ``A review and analysis of supply chain operations reference''~\cite{huan2004review} \hskip10em}\\
  \end{flushright}
  \hfill \hyperlink{scor_lv3}{\beamerbutton{Pg:SCOR Lv3}}

  %%%----
  \note{
    \textbf{Purpose:} Introduce to Level 2\\
    \textbf{Leading Question:} <<none>>\\
    \includegraphics [width=0.40\textwidth]{\figureLocation{scm_strategy/scor_cycleTime.jpg}}

    Level 1 defines scope of process, Lv2 defines config
    \begin{itemize}
      \item \textbf{MRO} maintenance-repair-operations
      \item \textbf{Flexility:} avg day per schedule change
    \end{itemize}
  }

\end{frame}

%%==============================================================================================================
\begin{frame}[label=scor_lv3]
  \frametitle{SCOR Level 3: Example and Insights \forIE}

  \NextHalfLine
  \begin{center}
    \only<1 |handout:1>{ \includegraphics [height=0.60\textheight]{\figureLocation{scm_strategy/scor_lv3_01.png}}}
    \only<2 |handout:0>{ \includegraphics [height=0.60\textheight]{\figureLocation{scm_strategy/scor_lv3_02.png}}}
    \only<3 |handout:0>{ \includegraphics [height=0.60\textheight]{\figureLocation{scm_strategy/scor_lv3_03.png}}}
    \only<4 |handout:0>{ \includegraphics [height=0.60\textheight]{\figureLocation{scm_strategy/scor_lv3_04.png}}}
  \end{center}

  \NextHalfLine
  %\textsc{OKC's Insights}
  \begin{flushleft} \footnotesize{
  \begin{itemize}
    \item \textbf{Plan:} requirements $\pm \approx \pm$  resources $\rightarrow$ plan %$\overbrace{\pm}$
    \item \textbf{Source$_{ETO}$ =} Source$_{MTO}$ + identify + select supplier
    \item \textbf{Delivery$_{ETO} \approx$ } Delivery$_{MTO}$ w/  negotiation +  sch. installation
    \item \textbf{Return$_{src}$ = } Return$_{dvy}$ + identify + disposition \alert{- transfer}
  \end{itemize} }
  \end{flushleft}

%
%
  %%%----
  \note{
    \textbf{Purpose:} Introduce to Level 2\\
    \textbf{Leading Question:} <<none>>\\
    \includegraphics [width=0.40\textwidth]{\figureLocation{scm_strategy/scor_cycleTime.jpg}}
    \begin{itemize}
      \item \textbf{MRO} maintenance-repair-operations
      \item \textbf{Flexility:} avg day per schedule change
    \end{itemize}
  }

\end{frame}

%=================================================================================================
\section[Strategy]{Supply Chain Strategy}
%=================================================================================================
\begin{frame}
  \frametitle{How a company compete with others?}

   \begin{itemize}
    \item \textbf{\code{Price}:} Wal-mart, Big Cola, Daiso
    \item \textbf{Agility \& Responsive:} Zara, Dell, JAL (on time)
    \item \textbf{Superb Quality:} Leica, Louis Vuitton
    \item \textbf{Innovation \& Flexibility:} Apple, Intel
    \item \textbf{Convenience \& Presence:} 7-Eleven, Starbucks, McDonald
    \item \textbf{Service \& Coverage:} 3M, AIS, Facebook, CitiBank, IG
   \end{itemize}

  \NextLine
  \begin{alertblock}{}{}
    What are strategic components to compete?
  \end{alertblock}

  \note{
    \textbf{Purpose:} Introduce to business strategy\\
    \textbf{Leading Question:} How the company compete\\
    \begin{itemize}
      \item \textbf{Note:} without differentiation (black bullets), customer decides by price
      \item \textbf{Flexibility}: change volume large (economy of scale $\leftrightarrow$ no layoff)/
      \item \textbf{more example:} Taiwan Semiconductor Manufacturing Company (TSMC), Samsung as OEM of Singer TV, India pharisaical company
    \end{itemize}

    \begin{tabular}{lrrr} \hline
      Strategy   & Advantage      & Compete                     & Key factor \\ \hline
      Innovation & brand \& tech  & desirable features/products & t2Market, t2Volumne \\
      Cost       & efficient oper & lowest prices               & cost/volumne \\
      Service    & superb service & tailor customer needs       & design from customer persp \\
      Quality    & safe \& reliable& no hassle                  & quality control \\ \hline
    \end{tabular}

  }
\end{frame}

%=================================================================================================
\begin{frame}[label=stra]
  \frametitle{Understanding strategic components }

   \begin{itemize}
    \item \textbf{Operation strategy:} how to produce goods and services?
    \item \textbf{Outsourcing strategy:} how to get products and services to buyers/ consumers?  \hfill \hyperlink{outs}{\beamerbutton{Pg:Outsourcing}}
    \item \textbf{Channel strategy:} what are in-house skills and expertise? \hfill \hyperlink{chan}{\beamerbutton{Pg:Channel \& Service}}
    \item \textbf{Customer service strategy:} what are customer's needs and how to turn into profits ? \hfill \hyperlink{chan}{\beamerbutton{Pg:Channel \& Service}}
    \item \textbf{Asset network:} what are organization of resources? \hfill \hyperlink{asset}{\beamerbutton{Pg:Asset}}
   \end{itemize}
  \begin{flushright}
    \tiny{\textbf{source:}  Cohen, S. and Roussel, J 2001. ``Strategic Supply Chain Management'' \hskip5em}\\
  \end{flushright}

  \note{
    \textbf{Purpose:} SCM strategy fitting with?\\
    \textbf{Leading Question:} how SCM strategy related to other strategy?\\
    \begin{itemize}
      \item \textbf{operation VS outsoucing:} make or buy?
      \item \textbf{rest} who and how to deliver 'value' and constraints
    \end{itemize}
  }
\end{frame}

%============================================================================================================
\begin{frame}
  \frametitle{Operation Strategy  }

  \begin{center}
  \begin{tikzpicture}[xscale=1.65,yscale=0.60,>=stealth]

    \draw[-, thick] (3.0,4.5) -- (3.0,-1.5) ;
    \draw[<->,thick] (-3.0,4.0) -- node[pos=0.5,auto, blue,sloped] {\textFont{ETO}} (3.0,4.0) ;
    \draw[<->,thick] (-1.8,3.0) -- node[pos=0.5,auto, blue,sloped] {\textFont{MTO}} (3.0,3.0) ;
    \draw[<->,thick] (-0.6,2.0) -- node[pos=0.5,auto, blue,sloped] {\textFont{ATO}} (3.0,2.0) ;
    \draw[<->,thick] (1.8,1.0) -- node[pos=0.5,auto, blue,sloped] {\textFont{MTS}} (3.0,1.0) ;

    \draw[<->,thick] (-3.0,-1.25) -- node[pos=0.5,auto, black,sloped] {\tiny{Engineering}} (-1.8,-1.25) ;
    \draw[<->,thick] (-1.8,-1.25) -- node[pos=0.5,auto, black,sloped] {\tiny{Procurement}} (-0.6,-1.25) ;
    \draw[<->,thick] (-0.6,-1.25) -- node[pos=0.5,auto, black,sloped] {\tiny{Fabrication}} (0.6,-1.25) ;
    \draw[<->,thick] (0.6,-1.25)  -- node[pos=0.5,auto, black,sloped] {\tiny{Assembly}} (1.8,-1.25) ;
    \draw[<->,thick] (1.8,-1.25)  -- node[pos=0.5,auto, black,sloped] {\tiny{Delivery}} (3.0,-1.25) ;
  \end{tikzpicture}
  \end{center}
  \begin{flushright}
    \tiny{\textbf{source:}Jacob and Chase 2010~\cite{jacobs2010operations}.}
  \end{flushright}


  \footnotesize{
  \begin{itemize}
    \item \textbf{ETO:} work with the customer to design and then make the product
    \item \textbf{MTO:} make the customer's product from raw mat'
    \item \textbf{ATO:} combine a number of preassembled modules to meet customer's specifications
    \item \textbf{MTS:} serve customers from finished goods inventory
  \end{itemize}
}

  \note{
    \textbf{Purpose:} Recap Operation strategy\\
    \textbf{Leading Question:} demand info $\rightarrow$ mfg process $\rightarrow$ delivery \\
    \begin{itemize}
    \item \textbf{ETO:} fit to specific location $\rightarrow$ product design, e.g., overhead crane
    \item \textbf{MTO:} expensive $\rightarrow$ part purchasing, e.g., super computer, airplane
    \item \textbf{ATO:} high variation $\rightarrow$ part customization, e.g., automobile, dell computer
    \item \textbf{MTS:} common goods/standardize $\rightarrow$ product delivery, e.g., TV, radio
    \end{itemize}
  }
  \end{frame}

%================================================================================================-
\begin{frame}
  \frametitle{Which operation strategy?}

  \begin{center}
    \includegraphics [height=0.45\textheight]{\figureLocation{scm_strategy/invPositbyInd.png}}
  \end{center}
  \begin{flushright}
    \tiny{\textbf{source:} Performance Measurement Group, 2004 \hskip10em }\\
  \end{flushright}

  \begin{center} \small{
    \begin{tabular}{lrr} \hline
      Type& Condition & Key benefits \\ \hline
      MTS & standard product w/ \alert{high volume}   & cost \& response \\
      ATO & ~~~~~~~~ product w/ \alert{high variation}& reduce inventory \\
      MTO & customized or infrequent demand product & low FG, many option \\
      ETO & complex product w/ \alert{unique customer} & response to requirement \\ \hline
    \end{tabular}}
  \end{center}

  \note{
    \textbf{Purpose:} diff business/ industry/ product $\rightarrow$ diff oper strategy\\
    \textbf{Leading Question:} <<none>>\\
    \begin{itemize}
      \item some industry has dominate oper. strategy, i.e., package food
      \item other industry is inconclusive
      \item good to group, but bad to decide
      \item \textbf{Note:} car (Faerie VS Honda motorcycle) and cooked foods (hi-end french VS military canteen) always bad example
    \end{itemize}
  }
\end{frame}

%=================================================================================================
\begin{frame}[label=outs]
  \frametitle{Outsourcing strategy \forIE}

  \begin{center} \small{
  \begin{tabular}{cc}\hline
    \textbf{benefit}                 & \textbf{concern} \\ \hline
        ramp-up/down (scope/scale)   & balancing value \& power  \\
        flexibility                  & culture \& execute \\
        focus on core                & proprietary \\ \hline
  \end{tabular} }
  \end{center}

	\NextHalfLine
  \begin{block}{}{}
    Every company eventually needs profit
  \end{block}

	\NextHalfLine
	\textsc{Should you outsource?}
  \begin{itemize}
    \item What give you competitive advantage (rolex, cisco)?
    \item What is your capability when fluctuation?
    \item What is your internal uniqueness of business process (BTV)?
    \item What is your power to negotiate (blink)? \hfill \hyperlink{stra}{\beamerbutton{Pg:Stratgic Components}}
  \end{itemize}

  \note{
    \textbf{Purpose:} Introduce to outsoucing\\
    \textbf{Leading Question:} should you outsource operation?\\
    \begin{itemize}
      \item how to reap benefit and prevent risk
      \item answer and discuss later in "soucing" lecture
    \end{itemize}
  }
\end{frame}

%=================================================================================================
\begin{frame}[label=chan]
  \frametitle{Customer Serv. \& Channel Strategy \forIE}

  \begin{center}
    \includegraphics [width=0.65\textwidth]{\figureLocation{scm_strategy/distChannel.png}}
  \end{center}

	\textsc{What are suitable channels?}
  \begin{itemize}
   \item Who is a targeted customer?
   \item What are costs and value of sale?
   \item What are channel partners should we have? \hfill \hyperlink{stra}{\beamerbutton{Pg:Stratgic Components}}
  \end{itemize}


  \note{
    \textbf{Purpose:} Introduce to channel\\
    \textbf{Leading Question:} <<none>>\\
    \begin{itemize}
      \item \textbf{direct} internet, tele-marketing, retailer
      \item \textbf{indirect} wholesale, reseller, field agents
      \item \textbf{level = numbers of parties:}
      \item \textbf{Agent:} extension of manufacturers can communicate with customer or other parties. \alert{no inventory}. Example of 4-Lv is Coke.
      \item \href{https://www.economicsdiscussion.net/distribution-channel/types-of-distribution-channels/31760}{\code{more info}}
    \end{itemize}

  }
\end{frame}

%=================================================================================================
\begin{frame}[label=asset]
  \frametitle{Asset Network \forIE}

  \begin{center}
    \includegraphics [width=0.75\textwidth]{\figureLocation{scm_strategy/autoSupplyChain.png}}
  \end{center}

  \NextHalfLine
  \textsc{Geography of Asset Network:} why multiple locations?
  \begin{itemize}
    \item \textbf{Global model} collaborate components , capital investment (intel)
    \item \textbf{Regional model:} trade agreement (Asian Auto)
    \item \textbf{Country model:} language, local content (TH Auto) \hfill \hyperlink{stra}{\beamerbutton{Pg:Stratgic Components}}
  \end{itemize}

  \note{
    \textbf{Purpose:} how and where to invest \\
    \textbf{Leading Question:} <<none>>\\
    \begin{itemize}
      \item \textbf{Balance:} Capital Intensive (cost of machine) \textsc{VS} Technology \textsc{VS} regulation \& trade block
    \end{itemize}
  }
\end{frame}

%=================================================================================================
\begin{frame}
  \frametitle{Strategic alignments }

  \begin{center}
    \includegraphics [width=0.70\textwidth]{\figureLocation{scm_strategy/alignment.png}}
  \end{center}

  \begin{itemize}
    \item \textbf{Business strategy:} how to win customers \& make money
    \item \textbf{Marketing strategy}: what is positioning of market and product
    \item \textbf{Supply chain strategy}: mean to fulfill customer's need
%    \begin{itemize}
%      \item \textbf{Operation:} operation strategy
%      \item \textbf{Procurement:} outsourcing, sourcing, partner
%      \item \textbf{Network:} asset, channel,
%    \end{itemize}
  \end{itemize}

  \note{
    \textbf{Purpose:} Business first. Then, market and supply chain together \\
    \textbf{Leading Question:} <<none>>\\
    \begin{itemize}
      \item \textbf{Business strategy:} how to survive in long term (as brain \& heart)
      \item \textbf{Mkt strategy}: how to engage customer (who? / story of company?) (as mouth and appearance)
      \item \textbf{SCM strategy}: how to manufacturing and delivery (as limbs)
    \end{itemize}
  }
\end{frame}

%=================================================================================================
\begin{frame}
  \frametitle{Alignment of SC and Business strategy }

%  \only<1 |handout:0>{
%    \begin{center}
%        \includegraphics [width=0.80\textwidth]{\figureLocation{scm_strategy/fiveForceExt.png}}
%    \end{center}
%   }
  \only<1 |handout:0>{
    \begin{center}
        \includegraphics [width=0.80\textwidth]{\figureLocation{scm_strategy/fiveForceInt.png}}
    \end{center}
   }
%  \only<3 |handout:0>{
%    \begin{center}
%        \includegraphics [width=0.80\textwidth]{\figureLocation{scm_strategy/SC_Enable_Framework.png}}
%    \end{center}
%   }
  \only<2 |handout:0>{
    \begin{center}
        \includegraphics [width=0.80\textwidth]{\figureLocation{scm_strategy/SCOR_cycle_Framework.png}}
    \end{center}
   }
  \only<3-|handout:1>{
  \begin{center} \footnotesize{
    \begin{tabular}{lrrrr} \hline \hline
    \textbf{Business} & \textbf{Source of} & \textbf{Basis of}    & \textbf{Key supply chain}  \\
    \textbf{Strategy} & \textbf{Advantage} & \textbf{Competition} & \textbf{Contribution}\\ \hline
    \textsc{Cost}     & cost-eff. operation& lowest cost          & efficient, cost improv  \\[2pt]
                      & \& mass            & \& economy of scale  &  risk \& global$^n$   \\[5pt]
    \textsc{Innovation}& brand, tech        & hype, desirable,     & time2market, volume  \\[2pt]
                      & \& niche           & \& innovative        &  \& manuf flex       \\[5pt]
    \textsc{Service}  & superb service     & fitted with need     & designed for customer\\[2pt]
                      &                    & \& after sale        & \& network           \\[5pt]
    \textsc{Quality}  & safest \& reliable & always work          & SC excellence, QC    \\[2pt]
                      &                    &  \& durable          & \& source            \\[5pt]
    \textsc{Agile}    & speed \& flex      & responsive           & locations, mode      \\[2pt]
                      &                    &                      & \& sync              \\[5pt]  \hline \hline
    \end{tabular} }
  \end{center}
  }
  \begin{flushright}
    \hfill \tiny{\textbf{source:} Johnson, M. and Stevens, G. ``Integrating the Supply Chain... 25 years on''  \hskip10ex~}
  \end{flushright}

  \note{
    \textbf{Purpose:} standard SCM and business strategy based on Porter 5 Forces\\
    \textbf{Leading Question:} what are classical SCM strategy?\\
    \begin{itemize}
      \item \textbf{Porter 5 Forces:} competition $\propto$ supplier, customer, substitution, barrier,internal
      \item \textbf{Innovative:} urge to buy a new innovation based on its brand $\rightarrow$ reliable and quick tech product
      \item[] \hskip3ex capital \& brand loyalty as  barrier \& substitution
      \item \textbf{Cost:} cheap to buy well-known product $\rightarrow$ low margin cost = $\frac{\mbox{value}}{\mbox{cost}}$
      \item[] \hskip3ex small margin to keep out as barrier \& substitution
      \item \textbf{Service:} comprehensive product/ service $\rightarrow$
      \item[] \hskip3ex back-online to keep out as barrier \& substitution
      \item \textbf{Quality:} lasting-reliable product; easy to fix $\rightarrow$ well-known product w/ second market
      \item[] \hskip3ex reputation \& record to keep out as barrier \& substitution
    \end{itemize}
  }
\end{frame}

%=================================================================================================
\begin{frame}
  \frametitle{Category and examples of SC Strategy}

  \begin{itemize}
    \item \textbf{Structure:} vertical integr$^n$, horizontal integr$^n$, virtual integr$^n$, collaborative,
    \item \textbf{Competitive Advantage:} low cost, differentiation
    \item \textbf{Market:} board, focus, best-value, platform, collaborative
    \item \textbf{Technology:} ecosystem, data centric, cyber-security, synchronization
    \item \textbf{Resilience:} multi-sourcing, near-shoring, harmonization, sustainability, buffer
    \item \textbf{Manufacturing:} diversification, globalization, localization, standard process
  \end{itemize}

  \NextLine
  \begin{alertblock}{}{}
    \begin{itemize}
      \item How to select to achieve \alert{Strategic Fit}? Examples \hfill \hyperlink{SC-Fit}{\beamerbutton{Pg:SC-Fit}}
      \item Is there any simplify or generic supply chain strategy? \hfill \hyperlink{SC-Strategy}{\beamerbutton{Pg:SC-Strategy}}
    \end{itemize}
  \end{alertblock}

  \note{
    \textbf{Purpose:} standard SCM and business strategy based on Porter 5 Forces\\
    \textbf{Leading Question:} what are classical SCM strategy?\\
    \begin{itemize}
      \item \textbf{Porter 5 Forces:} competition $\propto$ supplier, customer, substitution, barrier,internal
      \item \textbf{Innovative:} urge to buy a new innovation based on its brand $\rightarrow$ reliable and quick tech product
      \item[] \hskip3ex capital \& brand loyalty as  barrier \& substitution
      \item \textbf{Cost:} cheap to buy well-known product $\rightarrow$ low margin cost = $\frac{\mbox{value}}{\mbox{cost}}$
      \item[] \hskip3ex small margin to keep out as barrier \& substitution
      \item \textbf{Service:} comprehensive product/ service $\rightarrow$
      \item[] \hskip3ex back-online to keep out as barrier \& substitution
      \item \textbf{Quality:} lasting-reliable product; easy to fix $\rightarrow$ well-known product w/ second market
      \item[] \hskip3ex reputation \& record to keep out as barrier \& substitution
    \end{itemize}
  }
\end{frame}

%=================================================================================================
\section[SC strategy]{Strategic Fit and Supply Chain Strategy}
%=================================================================================================
\begin{frame}[label=SC-Fit]
  \frametitle{Achieving Strategic Fit \forIE}

   Alignment between \alert{business strategy} and \alert{supply chain}
   \begin{itemize}
    \item Understand \alert{nature of business}, i.e., customer, uncertainty, capacity
    \item Build \alert{sufficient capabilities},
    \item Deliver \alert{what customers value}
   \end{itemize}

  \NextLine
  \textsc{Examples}
  \begin{itemize}
    \item \textbf{Airline industry:} perishable capacity $\rightarrow$ dynamic price \& channels
    \item \textbf{Automotive industry:} variation, regulation, $\rightarrow$ JIT, module
    \item \textbf{Consumer electronic industry:} innovation, short life cycle $\rightarrow$ roll-out, leased capacity
    \item \textbf{Retail industry:} price \& availability $\rightarrow$ single DCs \& online, house brand
  \end{itemize}


  \note{
    \textbf{Purpose:} step for alignment\\
    \textbf{Leading Question:} how to translate these strategic into action \\
    \begin{itemize}
      \item \textbf{Airline:} time= expensive, must full as marginal cost
      \item \textbf{Auto:} limited variation; some local/some global;  $\exists$ critical components such as seat belt, engine block,
      \item \textbf{Consumer Electronic:} ; $\exists$ critical compony AMD \textsc{VS} Intel, AppMaterial
      \item \textbf{Retailer:} low margin, focus on logistics
    \end{itemize}
  }
\end{frame}

%=================================================================================================
\begin{frame}
  \frametitle{More Examples of Strategic Fit \forIE}


  \only<1 |handout:0>{
    \begin{center}
        \includegraphics [height=0.75\textheight]{\figureLocation{scm_strategy/Stra_locCostAirline.png}}
    \end{center}
   }
  \only<2 |handout:0>{
    \begin{center}
        \includegraphics [height=0.75\textheight]{\figureLocation{scm_strategy/Stra_ikea.png}}
    \end{center}
   }
  \only<3 |handout:1>{
    \begin{center}
        \includegraphics [height=0.75\textheight]{\figureLocation{scm_strategy/Stra_zara.png}}
    \end{center}
   }
  \begin{flushright}
    \hfill \tiny{\textbf{source:} Porter, M.  \hskip10ex~}
  \end{flushright}

  \note{
    \textbf{Purpose:} step for alignment\\
    \textbf{Leading Question:} how to translate these strategic into action \\
    \begin{itemize}
      \item \textbf{Airline:} time= expensive, must full as marginal cost
      \item \textbf{Auto:} limited variation; some local/some global;  $\exists$ critical components such as seat belt, engine block,
      \item \textbf{Consumer Electronic:} ; $\exists$ critical compony AMD \textsc{VS} Intel, AppMaterial
      \item \textbf{Retailer:} low margin, focus on logistics
    \end{itemize}
  }
\end{frame}

%=================================================================================================
\begin{frame}
  \frametitle{Obstacles to Strategic Fit \forIE}

  \begin{center}
    \only<1|handout:1>{\includegraphics [height=0.5\textheight]{\figureLocation{scm_strategy/StrategicFit.png}} }
  \end{center}
  \begin{flushright}
    \tiny{\textbf{source:} Chopra and Meindl 2010~\cite{chopra2007supply}\hskip10em}\\
  \end{flushright}


  \begin{itemize}
    \item \textbf{Business Nature:} variability (GM, FMCG), new feature (phone)
    \item \textbf{Fragmentation:} many ownership, too much outsourcing e.g., K-Mart
    \item \textbf{Globalization:} price, regulation, e.g., clothes
    \item \textbf{Business environment:} tech, e.g., internet, EV
  \end{itemize}

  \note{
    \textbf{Purpose:} dynamic/shifted ground also applies to stratgy\\
    \textbf{Leading Question:} Why some company fail on their SC strategy? \\
    \begin{itemize}
      \item \textbf{more SKU:} need more customer $\rightarrow$ complex
      \item \textbf{short LC:} $\rightarrow$ little time to perfect
      \item \textbf{more competitive:} cheaper, faster, better, safer, convince
      \item \textbf{less risk:} share risk/more party
      \item \textbf{more info/play field}
    \end{itemize}
  }
\end{frame}

%=================================================================================================
\begin{frame}[label=SC-Strategy]
  \frametitle{Understanding demand and supply}

  \NextLine
  \begin{center}
  \begin{tabular}{lcc}
    Demand  &  \includegraphics [width=0.15\textwidth]{\figureLocation{scm_strategy/whiteShirt.png}}
            &  \includegraphics [width=0.15\textwidth]{\figureLocation{scm_strategy/zaraShirt.png}}\\
            &  \tiny{functional products} & \tiny{innovative products} \\
    Supply  & \only<1|handout:1>{\includegraphics [width=0.25\textwidth]{\figureLocation{scm_strategy/hydroPlant.png}}}
              \only<2|handout:0>{\includegraphics [width=0.25\textwidth]{\figureLocation{scm_strategy/sewingProcess.png}}}
            & \only<1|handout:1>{\includegraphics [width=0.25\textwidth]{\figureLocation{scm_strategy/gasPlant.png}} }
              \only<2|handout:0>{\includegraphics [width=0.25\textwidth]{\figureLocation{scm_strategy/compChipProcess.jpg}}}\\
            & \tiny{stable processes} & \tiny{evolving processes} \\
  \end{tabular}
  \end{center}

  \begin{block}{\footnotesize{Question 4}} \footnotesize{
  \begin{itemize}
    \item What are characteristics of \alert{functional products} and \alert{innovative products }
    \item What are characteristics of \alert{stable processes} and \alert{evolving processes}
  \end{itemize} }
  \end{block}


  \note{
    \textbf{Purpose:} Introduce to classical supply chain strategy\\
    \textbf{Leading Question:} how demand/supply affect supply chain and its strategy?\\
    \begin{itemize}
      \item \textbf{demand uncertain:} high functional/R\&D $\rightarrow$ innovative
      \item \textbf{supply uncertain:} evolving process $\rightarrow$ low standard
    \end{itemize}
  }
\end{frame}

%=================================================================================================
\begin{frame}
  \frametitle{Uncertainty Matrix}

  \begin{center}
    \begin{tabular}{ll|p{0.35\textwidth}|p{0.35\textwidth}|}
      \multirow{6}{*}{\rotatebox{90}{\textbf{Supply Processes}} }
                        & \multicolumn{1}{c}{}
                        & \multicolumn{2}{c}{\textbf{Demand Fluctuation}} \\
                        &  \multicolumn{1}{c}{}
                        & \multicolumn{1}{c}{\textsl{Functional}}
                        & \multicolumn{1}{c}{\textsl{Innovative}} \\ \cline{3-4}
      & \textsl{Stable} &  oil, gas, food      & fashion apparel,  \\
      &                 &  basic clothes       & computer, music  \\[5pt]  \cline{3-4}
      & \textsl{Evolve} &  hydro electric,     & telecommunication \\
      &                 &  some food           & iPhone, iPad, CPU \\[5pt]  \cline{3-4}
    \end{tabular}
  \end{center}
  \begin{flushright}
    \tiny{\textbf{source:} ``Aligning Supply Chain Strategies with Product Uncertainties''. \emph{Lee, H.L.} 2002~\cite{lee2002aligning} \hskip10em}\\
    \tiny{\textbf{\hskip5ex} ``Introduction to Supply Chain Management''. \emph{Raz, Gal} 2009 \hskip10em}\\
  \end{flushright}

  \note{
    \textbf{Purpose:} Introduce to products based on demand/supply uncertainty\\
    \textbf{Leading Question:} How do we apply the concept? Why product?\\
    \begin{itemize}
      \item \textbf{Oil:} constant demand - proven tech
      \item \textbf{Hydro:} constant demand - seasonal suuply
      \item \textbf{Fashion:} shifted style - 100yr old tech
      \item \textbf{iPhone:} saturated/need replaced smart phone-change techn 300nm last 20 years $\rightarrow$ 7nm
    \end{itemize}
  }
\end{frame}

%=================================================================================================

\begin{frame}
  \frametitle{Supply chain capability}

  \begin{center}
    \includegraphics [width=0.8\textwidth]{\figureLocation{scm_strategy/ResponsiveSpectrum.png}}
  \end{center}
  \begin{flushright}
    \tiny{\textbf{source:} Chopra and Meindl 2010~\cite{chopra2007supply} \hskip10em}\\
  \end{flushright}


  \begin{itemize}
    \item \textbf{Efficient/Lean:} lean \& cost driven, e.g. grocery
    \item \textbf{Responsive:} flexible and tailor to needs of customer, e.g. fashion apparel
    \item \textbf{Risk-Hedging:} pooling and sharing resources with alternative sources, e.g. hydro-power
    \item \textbf{Agile:} combining responsive and risk hedging, e.g. high tech, telecom (react \& quick execute)
  \end{itemize}

  \note{
    \textbf{Purpose:} Spectrum of SC strategy\\
    \textbf{Leading Question:} what are four SC strategy?\\
    \begin{itemize}
      \item $<<$explain bullets $>>$
    \end{itemize}
  }
\end{frame}

%=================================================================================================
\begin{frame}
  \frametitle{Lean SC VS Agile SC}

  \begin{center}
    \begin{tabular}{cc}
      \includegraphics [width=0.25\textwidth]{\figureLocation{scm_strategy/leanSC.png}}~~
    &
      ~~\includegraphics [width=0.25\textwidth]{\figureLocation{scm_strategy/agileSC.png}}    \\
    \end{tabular}
  \end{center}
  \begin{flushright}
    \tiny{\textbf{source:} Chopra and Meindl 2010~\cite{chopra2007supply}\hskip10em}\\
  \end{flushright}

  \begin{center}
  \small{
  \begin{tabular}{lrr} \hline
                  & Lean SC   & Agile SC \\  \hline
    product       & community & fashion \\
    cust.driver   & cost      & other \\
    profit margin & low       & high \\
    main cost     & operation & marketing \\ \hline
  \end{tabular}
  }
  \end{center}

  \note{
    \textbf{Purpose:} Introduce two dominated/well-known strategy \\
    \textbf{Leading Question:} What are typical/frequent strategy? \\
    \begin{itemize}
      \item \textbf{Lean SC Strategy:} high volume \& low variety/ planning consistency
      \item \textbf{Agile SC Strategy:} high variety \& low volume/ quick reactive
      \item \textbf{Note I:} lean $\neq$ agile $\rightarrow$ lagile (research topic)
      \item \textbf{Note II:}
      \begin{itemize} \scriptsize{
        \item There is a cost to achieving responsiveness
        \item Supply chain efficiency: cost of making and delivering the product to the customer
        \item \alert{Increasing responsiveness} results in \alert{higher costs} that lower efficiency }
      \end{itemize}
    \end{itemize}
  }
\end{frame}

%%=================================================================================================
\begin{frame}
  \frametitle{Principles of Lean Supply Chain}

  \begin{center}
    \only<1>{ \includegraphics [height=0.5\textheight]{\figureLocation{scm_strategy/7wastes.png}} }
    \only<2|handout:0>{\includegraphics [height=0.5\textheight]{\figureLocation{scm_strategy/valueStreamMapping.jpg}}}
  \end{center}

  \begin{flushright}
    \tiny{\textbf{source:} Chopra and Meindl 2010~\cite{chopra2007supply}\hskip10em}\\
  \end{flushright}

  \small{
  \begin{itemize}
    \item \textbf{Principle:} eliminating wastes by marginal gain, bottom up, training
    \item \textbf{7Wastes:} TIM WOOD B.
    \item \textbf{Basic:} 5S, layout, work procedure with visual,
    \item \textbf{Adv.} value stream mapping, reducing batch size \& \# suppliers
    \item \textbf{Logistics:} crossdock, back haul, postponement
  \end{itemize} }

  \note{
    \textbf{Purpose:} Introduce to lean supply chain \\
    \textbf{Leading Question:} <<none>>\\
    \begin{itemize}
      \item \textbf{history of lean SCM:} toyota $\rightarrow$ office $\rightarrow$ process
      \item \textbf{7Wastes:} defects, production, transportation, inventory, motion, process, waiting
      \item \textbf{SC Tech}: JIT delivery, kitting for POU, VMI, 3PL
      \item \alert{Increasing responsiveness} results in \alert{higher costs} that lower efficiency
      \item \textbf{Scope:} add \alert{value} (i.e., physical changes, customer cares)
      \begin{itemize} \scriptsize{
        \item focus on \alert{people/department} who directly add value
        \item involve suppliers in R\&D and logistics
        \item optimize across organization}
      \end{itemize}
    \end{itemize}
  }
\end{frame}

%%=================================================================================================
\begin{frame}
  \frametitle{Principles of Agile Supply Chain}

  \begin{center}
    \only<1 |handout:1>{\includegraphics[height=0.6\textheight]{\figureLocation{scm_strategy/agileSC_concept02.png}}}
    \only<2-|handout:0>{\includegraphics[height=0.6\textheight]{\figureLocation{scm_strategy/agileSC_integrate01.png}}}
  \end{center}

  \begin{flushright}
    \tiny{\textbf{source:} Chopra and Meindl 2010~\cite{goetschalckx2011supply} \hskip10em}\\
  \end{flushright}

  \begin{itemize}
    \item \textbf{What:} rapidly response to \alert{unpredictable market demand} driven by innovation
    \item \textbf{Scope:} extend \alert{end-to-end} of supply chain
  \end{itemize}

  \note{
    \tiny{Agile SC framework: configuration and structure are not fixed (fast) \& for innovative products (smart)} \\
    \textbf{Purpose:} Introduce to agile supply chain\\
    \textbf{Leading Question:} <<none>>\\
    \begin{itemize}
      \item \textbf{virtual integration:} shared info, collaborated plan, end2end visibility
      \item \textbf{market sensitive:} know consumer, daily POS, capture emerging trend
      \item \textbf{process alignment:} internal process (forecasting $\rightarrow$ delivery)// market based forecast
      \item \textbf{dynamic:} capabilities for responsive / anticipate unpredictable change // short term
    \end{itemize}
  }
\end{frame}

%================================================================================================-
\section[Metrics]{Supply Chain Performance Matric}
%================================================================================================-
\begin{frame}
  \frametitle{SCM: Strategic Cost Management}

  \begin{quotation}
    If you can not measure it, you can not improve it.
  \end{quotation}
  \begin{flushright}
    \hfill \tiny{\textbf{source:} William Thomson (Lord Kelvin) \hskip10ex~}
  \end{flushright}

  \NextHalfLine
  %\textsc{Major SCM Cost Element}
  \begin{center}
    \only<1|handout:1>{\includegraphics [height=0.5\textheight]{\figureLocation{scm_strategy/business_supplychain.png}}}
    \only<2|handout:0>{\includegraphics [height=0.5\textheight]{\figureLocation{scm_strategy/perfFramework.png}}}
    \only<3|handout:0>{\includegraphics [height=0.5\textheight]{\figureLocation{scm_strategy/majorSCMCost.png}}}
  \end{center}

  \note{
    \textbf{Purpose:} To achieve better SCM, must measure\\
    \textbf{Leading Question:} <<none>>\\
    \begin{itemize}
      \item \textbf{matrix hierarchy:} linage from strategy $\rightarrow$ objective $\rightarrow$ KPI
      \item \textbf{perspective}: operation, financial, customer
      \item \textbf{Note}
      \begin{itemize} \scriptsize{
        \item 47\% implements SCM for cost reduction
        \item where can find cost? \#transaction or impact or \#touch or \# party}
      \end{itemize}
    \end{itemize}
  }
\end{frame}

%================================================================================================-
\begin{frame}
  \frametitle{Performance Matrix }

  \textsc{Perspective}
  \begin{itemize}
    \item \textbf{Customer Service:} indicator of ability to meet \alert{customers} timely needs and create values
    \item \textbf{Operational:} indicator of \alert{processes} and \alert{decisions} to meet customer expectations
    \item \textbf{Financial:} indicator of value creation for \alert{shareholders}
  \end{itemize}

  \NextHalfLine
  \textsc{Dimensions}
    \begin{itemize}
      \item \textbf{Cost:} \alert{capital} required to meet expected process and operation
      \item \textbf{Responsiveness /Time:} speed in which expectations is met
      \item \textbf{Reliability:} trust/\alert{quality} of service and accuracy of information
      \item \textbf{Agility:}  ability and flexibility to \alert{respond to external} influences
      \item \textbf{Asset Management Efficiency:} ability to efficiently utilize \alert{assets}
    \end{itemize}

  %%%------
  \note{
    \textbf{Purpose:} Introduce to matrix\\
    \textbf{Leading Question:} <<none>>\\
    \begin{itemize}
      \item \textbf{Perspective} from strategy
      \item \textbf{Dimension} from SCOR model
    \end{itemize}
  }
\end{frame}

%================================================================================================-
\begin{frame}
  \frametitle{Importance Performance Measure }

  \NextHalfLine
  Discuss the following popular SCM indicators in terms of perspective and dimension\\
  \begin{block}{}{}{}
    \begin{tabular}{ll}
      $\bullet$ inventory day                       & $\bullet$ forecasting accuracy   \\
      $\bullet$ cash-to-cash cycle time             & $\bullet$ overall value at risk      \\
      $\bullet$ on-time delivery                    & $\bullet$ order fulfillment lead time \\
      $\bullet$ return of equity (ROE)              & $\bullet$ rate of good return \\
      $\bullet$ perfect order fulfillment \hskip5ex & $\bullet$ $\frac{\mbox{logistic cost}}{\mbox{cost of good sold}}$ \\
    \end{tabular}
  \end{block}

  Elaborate activities or department effected by these indicators
  \begin{alertblock}{\small{Question 5}}{}
    \begin{itemize}
      \item Elaborate activities or department effected by these indicators
      %\item Does Thailand have any benchmark?     \hfill \hyperlink{benchmark}{\beamerbutton{Pg:Benchmark}}
      \item How should we improve SC of our company?  \hfill \hyperlink{improve}{\beamerbutton{Pg:Design Improve}}
    \end{itemize}
  \end{alertblock}

  %%%------
  \note{
    \textbf{Purpose:} example of popular KPIs/ research these KPIs\\
    \textbf{Leading Question:} what are popular matrix? \\
    \begin{itemize}
      \item \textbf{Note:} always ask how to calculate
      \item \textbf{Reliability:} perfect order fulfillment
      \item \textbf{AM}: cash2cash cycle time, Return on supply chain fixed assets
    \end{itemize}
  }
\end{frame}

%================================================================================================-
\begin{frame}
  \frametitle{Benefits of Performance Measure }

   \begin{itemize}
    \item \textbf{Learning:} comparing and evaluating an organization \alert{internally}
    \item \textbf{Coordination:} indicating 'weak' points and \alert{re-allocate resources}
    \item \textbf{Motivation \& Goal:} communicating business strategy to operational \alert{goal}
    \item \textbf{Benchmark:} comparing 'our' performance and adopting \alert{best practice}
  \end{itemize}

  \NextLine
  \textsc{OKC's Note on Measurement}
  \begin{itemize}
    \item \textbf{Measuring = Non-Value Added:} require \alert{time}, effort, \alert{tool}, revision
    \item \textbf{Measuring $\rightarrow$ Action:} aware of top-down (econ$_{K9}$, art), find \alert{conflicts}
    \item \textbf{RE-Method in Business:} KPI \textsc{VS} Balance Scorecard \textsc{VS} OKR
  \end{itemize}

  %%%------
  \note{
    \textbf{Purpose:} Introduce to textbook\\
    \textbf{Leading Question:} <<none>>\\
    \begin{itemize}
      \item \textbf{Learning:} historical tracking/ possible for analysis and trouble shooting
      \item \textbf{Coordination \& Motivation \& Goal:} similar to budgeting, internal measuring in case conflict
      \item \textbf{Benchmark:} comparing company to the best in other area/ same area (Glillet/DHL country party Story)
      \item \textbf{Objective Key Result:} pick 3-5 top key numerical result / popularized by google/ 70\% = business as usual
      \item \textbf{Balance Scorecard:} multi-facet strategic perf measurement comprise of Financial, Customer, Operation, Growth
    \end{itemize}
  }
\end{frame}

%%================================================================================================-
%\begin{frame}
%  \frametitle{Office of Logistics ILPI \forIE}
%
%  \begin{center}
%  \begin{tabular}{cc}
%    \multicolumn{2}{c}{\only<1|handout:0>{\includegraphics [width=0.8\textwidth]{\figureLocation{scm_strategy/ilpi01.png}}}} \\
%    \only<2|handout:1>{\includegraphics [width=0.45\textwidth]{\figureLocation{scm_strategy/ilpi02.png}}} &
%    \only<2|handout:1>{\includegraphics [width=0.45\textwidth]{\figureLocation{scm_strategy/ilpi03.png}}}
%  \end{tabular}
%  \end{center}
%
%%  \begin{flushright}
%%   % \hfill \tiny{ \textbf{source:} \href{https://dol.dip.go.th/uploadcontent/DOL/Pert/ILPI\%202019\%20.pdf}{\code{https://dol.dip.go.th}} \hskip10ex~}
%%  \end{flushright}
%
%  %%%------
%  \note{
%    \textbf{Purpose:} Introduce to textbook\\
%    \textbf{Leading Question:} <<none>>\\
%    \begin{itemize}
%      \item \textbf{Story:} dept of logistic groups and trys to benchmark industry in thailand
%      \item \textbf{LPI} Industrial Logistic Performance Index /
%    \end{itemize}
%  }
%\end{frame}

%================================================================================================-
\section[Bullwhip Effect]{Information and Bullwhip Effect}
%================================================================================================-
\begin{frame}
  \frametitle{Bullwhip Effect}

  \begin{center}
    \only<1|handout:0>{\includegraphics [height=0.45\textheight]{\figureLocation{scm_strategy/distChannel01.png}}}
    \only<2|handout:0>{\includegraphics [height=0.45\textheight]{\figureLocation{scm_strategy/OKCbeergame.png}}}
    \only<3->{
      \begin{itemize}
        \item \textbf{What:} \alert{distortion} of \alert{demand information} within SC
        \item \textbf{Nature:} distortion \alert{increases in} higher \alert{upstream} of SC
        \item \textbf{Causes:}
        \begin{itemize}
          \item Information delay: \alert{uncertainty} in quantities and pattern of demand
          \item Length of chain: \alert{aggregated demands} of each party
        \end{itemize}
        \item \textbf{Effects:}  inventory (dead stock, backlog), production smoothness, expedition
        \item \textbf{Examples:} Proctor \& Gamble (Pampers); HP (printers); Barilla (pasta)
        \item \textbf{Remedy:} coordination (POS, forecasting), shorten SC, volume discount in fixed period
      \end{itemize}
    }
  \end{center}

  %%%------
  \note{
    \textbf{Purpose:} Introduce to MIT beer game/ a classical example of bull whipeffect\\
    \textbf{Leading Question:} <<none>>\\
    \begin{itemize}
      \item Increasing overall inventory $\rightarrow$ \alert{dead stock}
      \item Increasing transportation cost/ production
      \item Decreasing response time
      \item Evolving into a game between parties in supply chain
      \item \textbf{course:} promotion, price change
    \end{itemize}
  }
\end{frame}

%================================================================================================-
\begin{frame}
  \frametitle{Demand Fluctuation}

  \RemoveLine
  \begin{center}
    \includegraphics[width=0.75\linewidth]{\figureLocation{scm_strategy/BullWhipEffect.png}}\\
  \end{center}
  \begin{flushright}
    \tiny{\textbf{source:} Chase and Jacob. 2011.\hskip10em}\\
  \end{flushright}

  %%%------
   \note{
    \textbf{Purpose:} Show inventory result of game\\
    \textbf{Leading Question:} <<none>>\\
    \begin{itemize}
      \item <<explain figures>>
    \end{itemize}
  }
\end{frame}

%=================================================================================================
\begin{frame}[label=improve]
  \frametitle{Art of Supply Chain Improvement }

    \textsc{Supply Chain Basic Steps}
    \begin{itemize}
      \item \textbf{[D] Begin with customers:} need of \code{every customers} (Amazon.com, Zippo.com)
      \item \textbf{[D] Keep SCM simple:} focus on process (i.e., purchase, analyze, transport), not function
      \item \textbf{[I] Measure what important, and improve them:}
      \item \textbf{[I] Build bridge, eliminate wall:}  Realtime communication (UF), quantify investment, focus on people
    \end{itemize}

    \textsc{Supply Chain Advance}
    \begin{itemize}
      \item \textbf{[D] Focus on core, outsource others:}
      \item \textbf{[D] Aware roles change:} technology, opportunity, new comer
      \item \textbf{[D] Catch the vision, map it out:} know critical weak spot/ people, what value created
      \item \textbf{[I] Use Tech. as enabler:}
      \item \textbf{[I] Remind Fair $\neq$ equal relationship:}
    \end{itemize}
  \begin{flushright}
    \tiny{\textbf{source:} Fawcett, S. and Magnan, G. 2004.\hskip10em~}\\
  \end{flushright}

  %%%------
  \note{
    \textbf{Purpose:} Step to improve SCM: Basic/Advance\\
    \textbf{Leading Question:} How you would improve any SCM? \\
    \begin{itemize}
      \item \textbf{D} for design SCM, \textbf{I} for implement
      \item \textbf{Note:} in practice, there no clear separation between Basic and Advance, and may repeat
      \item human is main SCM problem; \#SKU $\rightarrow$ more complex \& difficult \& IT
      \item $<<$explain bullets$>>$
    \end{itemize}
  }
\end{frame}

%================================================================================================-

\begin{frame}
  \frametitle{Reference}

  \begin{footnotesize}
    \bibliography{\refLocation{Coordination},\refLocation{logisticProvider}}
  \end{footnotesize}

  \note{
    [empty] note for citation page
  }

\end{frame}

%=================================================================================================
\end{document}















%=================================================================================================
\begin{frame}
  \frametitle{Frame Title Here}

  \note{
    \begin{itemize}
      \item info
    \end{itemize}
  }
\end{frame}

%=================================================================================================
\begin{frame}
  \frametitle{Responsibility}

    \begin{center}
    \begin{tabular}{rr} \hline
      7-Eleven   & Owner   \\ \hline
      \alert{Supply products}
    &
      \alert{Operate} \& \alert{maintain} store
    \\
      Accounting, Advertising, IT Systems
    &
      Hiring staff \& customer service
    \\
      Install and remodel
    &
      Order products
    \\ \hline
    \end{tabular}
    \end{center}

  \note{
    \begin{itemize}
      \item info
    \end{itemize}
  }
\end{frame}

%=================================================================================================


%=================================================================================================
\begin{frame}
  \frametitle{Efficient and Responsive in Supply Chains}

  \begin{center}
  \small{
  \begin{tabular}{lP{0.25\textwidth} P{0.25\textwidth}P{0.25\textwidth} } \hline
    	                 & 	Efficient	           &  Responsive      & Differentiation \\ \hline
    ~~~Primary goal	   &  cost	               &  speed           & brand \& R\&D \\
    Design 	           &  min prod. cost	     &  short setup     & modularity\\
    Pricing	           &  low margin	         &  high margin     & high margin\\
    Manufacturing 	   &  high util.	         &  flex. capacity  & mass customization \\
    Inventory 	       &  min inventory	       &  buffer          & avoid obsolete  \\
    Lead time 	       &  N/A	                 &  short manuf.    & short R\&D\\
    Selection          &  cost $\rightarrow$ quality	
                       &  speed, flex, quality
                       &  joint partner \\
    Transportation 	   &  low-cost mode	       &  fast modes      & depend \\ \hline
  \end{tabular}
  }
  \end{center}

  \note{
    \begin{itemize}
      \item info
    \end{itemize}
  }
\end{frame}



%=================================================================================================
\begin{frame}
  \frametitle{Seafood supply chain}

  \RemoveLine
  \begin{center}
    \includegraphics [width=0.90\textwidth]{\figureLocation{seafoodSC.png}}
  \end{center}
  \begin{flushright}
    \tiny{\textbf{source:} Sal Forest Report \hskip10em }\\
  \end{flushright}

  \note{
    \begin{itemize}
      \item important and motivation of SCM~\cite{jacobs2010operations}~\cite{ballou2007evolution} \cite{chopra2007supply} \cite{goetschalckx2011supply}
      \item SCM is an elective class/key application of ICE
    \end{itemize}
  }
\end{frame}
%=================================================================================================

\begin{frame}
  \frametitle{Petrochemical supply chain}

  \RemoveLine
  \begin{center}
    \includegraphics [width=0.90\textwidth]{\figureLocation{oilSC.png}}
  \end{center}
  \begin{flushright}
    \tiny{\textbf{source:} Oil Industry Overview \hskip10em }\\
  \end{flushright}

  \note{
    \begin{itemize}
      \item important and motivation of SCM~\cite{jacobs2010operations}~\cite{ballou2007evolution} \cite{chopra2007supply} \cite{goetschalckx2011supply}
      \item SCM is an elective class/key application of ICE
    \end{itemize}
  }
\end{frame}

%=================================================================================================
\begin{frame}
  \frametitle{How to archived strategic fit? }

  \textsc{1. Understanding customer and supply chain uncertainty:} \alert{nature of business}
  \begin{alertblock}{}{}
    Attribute of customer demand + demand uncertainty $\rightarrow$ supply chain uncertainty
  \end{alertblock}

  \NextHalfLine
  \textsc{2. Understanding supply chain capability:} \alert{how a company meet the nature}
  \begin{alertblock}{}{}
    Responsiveness \textbf{VS} Efficiency \textbf{VS} Differentiate
  \end{alertblock}

  \NextHalfLine
  \textsc{3. Archiving Strategic fit:} \alert{match uncertainty with capability}

  \note{
    \begin{itemize}
      \item Identify \alert{needs} and \alert{customer segment}
      \item Quantify \alert{lot size}
      \item Sniff 'acceptable' \alert{response time} and \alert{service level}
      \item Determine suitable \alert{product variety} and \alert{innovated cycle}
      \item Select/Accept \alert{price} of product
    \end{itemize}
  }
\end{frame}




